\usepackage{booktabs}
\usepackage{float}
\usepackage{afterpage}

% Suppress the default pandoc/LaTeX title page. pandoc emits \maketitle from the
% YAML title/author/date, which produces a redundant title page (stamped with the
% build date) in front of our own front matter. The real title, copyright, and
% dedication pages are supplied by frontpage.tex (generated from index.rmd and
% included via before_body), so make \maketitle a no-op here.
\renewcommand{\maketitle}{}

% Advance to the next recto (odd) page, leaving any skipped verso completely
% blank -- no header, no page number. Used in the generated front matter
% (frontpage.tex) to put the title page and dedication on rectos, with the
% copyright page on the verso of the title page.
\newcommand{\clearrecto}{\clearpage\ifodd\value{page}\else\thispagestyle{empty}\null\clearpage\fi}

% Define book size
% margin sets the base 1in margin on all sides; bindingoffset adds a 0.25in
% gutter to the inner (spine) edge, mirrored on recto/verso, so the inside
% margin is 1.25in and the outside 1in -- the conventional allowance for
% perfect binding.
\usepackage[paperwidth=7in, paperheight=10in, margin=1in, bindingoffset=0.25in]{geometry}

% Running headers.
% The default book class typesets running heads in ALL CAPS at full size. On
% this narrow (~5in) text block, long section titles then run into the page
% number (e.g. "...DISCRETE CHARACTERS101"). Use smaller, mixed-case running
% heads so long titles fit, with the page number on the outer edge and the
% running title on the inner edge (further from the page number).
\usepackage{fancyhdr}
\setlength{\headheight}{14pt}
\pagestyle{fancy}
\fancyhf{}
\fancyhead[LE,RO]{\small\thepage}
\fancyhead[RE]{\small\nouppercase{\leftmark}}
\fancyhead[LO]{\small\nouppercase{\rightmark}}
\renewcommand{\headrulewidth}{0.4pt}
% Chapter-opening pages use the "plain" style; put their page number in the foot.
\fancypagestyle{plain}{%
  \fancyhf{}%
  \fancyfoot[C]{\small\thepage}%
  \renewcommand{\headrulewidth}{0pt}%
}

% Widow and orphan control.
% A "widow" is the last line of a paragraph stranded at the top of a page; an
% "orphan" (club line) is the first line of a paragraph stranded at the bottom.
% 10000 is TeX's "infinite" penalty, telling it to avoid such breaks wherever it
% can. We also guard the lines around displayed equations, of which this book has
% many.
\widowpenalty=10000
\clubpenalty=10000
\displaywidowpenalty=10000
\postdisplaypenalty=10000

% Keep section headings from being stranded near the foot of a page with only a
% line or two of their text following. \clubpenalties makes the penalty apply to
% the first *two* lines after a break (i.e. after a heading), not just the first,
% so a heading pulls at least a couple of lines with it to the next page.
\clubpenalties 3 10000 10000 0

% The book class uses \flushbottom (every page flush to the bottom margin), which
% leaves TeX no vertical slack to honour the penalties above -- it will still
% break rather than under-fill a page. \raggedbottom lets pages end a line or two
% short so the penalties can actually take effect. The trade-off is that facing
% pages may not always end at exactly the same height.
\raggedbottom